\chapter{Complexity — Step Block Forward}
%%% generated by the blueprint pipeline from TCSlib.Complexity.TuringMachine.Build.StepBlockForward
\section{Overview}

Declarations of the Lean module \texttt{TCSlib.Complexity.TuringMachine.Build.StepBlockForward}.

\section{Declarations}

\begin{definition}[Forward-loop initial configuration]
\lean{Turing.StepBlock.forwardLoopCfg}
\label{Turing.StepBlock.forwardLoopCfg}
\leanok
\uses{Turing.Cfg, Turing.FinTM.bufferTape, Turing.MultiTapeTM, Turing.StepBlock.summaryCells, Turing.tableBitsK, Turing.unaryFin}
Let $k,n\in\mathbb{N}$, let $Q$ be a state type, let $x$ be a Boolean input word, let
$M$ be a $k$-tape machine with Boolean work symbols and states in $\text{Fin}(n+1)$, and
let $c$ be a configuration of that machine on $x$. For an entry state $q\in Q$, the
forward-loop configuration has state $q$ when $c$ has a state and is otherwise halted;
it preserves the input-head position, the original work tapes and their head positions,
and the output. It adds two work tapes with head positions zero: one containing the
designated encoding of $M$'s transition table, and one containing the designated summary
initialized with no output and the unary encoding of $c$'s state when present.
\end{definition}

\begin{theorem}[A forward-loop simulation with bounded overhead]
\lean{Turing.StepBlock.exists_forwardLoopTM}
\label{Turing.StepBlock.exists_forwardLoopTM}
\leanok
\uses{Turing.Cfg, Turing.Lockstep, Turing.MultiTapeTM, Turing.MultiTapeTM.runFrom, Turing.MultiTapeTM.visitedByTapeHead, Turing.StepBlock.forwardLoopCfg, Turing.tableBitsK}
For every natural number $k$, there exist a finite state set $Q$ with decidable
equality, a Turing machine $L$ with $k+2$ work tapes, an entry state in $Q$, and a
natural number $C$ with the following property. For every natural number $n$, Boolean
input word $x$, and $k$-tape machine $M$ whose states are indexed by
$\mathrm{Fin}(n+1)$, there is a natural-number cost function on configurations such that
the simulation initialized by the forward-loop configuration map advances from the image
of each nonhalting configuration to the image of its successor under $M$ in exactly that
configuration's cost, maps halting configurations to halting configurations, and has no
halting intermediate state during any such simulated block. The cost of each nonhalting
configuration is positive and at most $50(|T|+1)+C$, where $T$ is the transition-table
serialization of $M$. Moreover, for every configuration $c$ and time $t$ during which
$M$ has not halted at any earlier time, each original tape head visits only cells
visited by the corresponding head of $M$ through time $t$; each of the two additional
tape heads visits only cells in the integer interval from $-(50(|T|+1)+C)$ to
$50(|T|+1)+C$ by the simulated time equal to the sum of the first $t$ costs along the
run of $M$ from $c$.
\end{theorem}
